\section{Conclusiones y Trabajos Futuros}

\subsection{6.1 Conclusiones Principales}

Resulta particularmente llamativo comprobar, tras completar el ciclo completo de diseño, simulación y análisis, hasta qué punto la fabricación aditiva puede redefinir los límites de lo posible en plataformas satelitales de telecomunicaciones. Los resultados obtenidos validan de forma consistente la hipótesis central: una aproximación sistemática de Design for Additive Manufacturing permite reducciones de masa del orden del 27-32\% mientras se preservan —e incluso mejoran— las prestaciones estructurales, dinámicas y RF.

Boschetto et al. \cite{Boschetto2023} ya habían señalado el potencial de las estructuras lattice con thickening estratégico; este trabajo confirma y extiende esos hallazgos al contexto específico de minisatélites para SATCOM. La consolidación de múltiples piezas en componentes multifuncionales no solo aligera la estructura, sino que reduce drásticamente interfaces críticas, tiempos de ensamblaje y potenciales puntos de fallo. Los diagramas de ganancia y el ancho de banda logrado en bandas Ku demuestran que la integración de requisitos multidisciplinarios ya no representa un compromiso inevitable, sino una oportunidad real de optimización global.

En términos cuantitativos, la metodología propuesta entrega un balance favorable entre desempeño técnico y viabilidad económica, especialmente en escenarios de producción en lotes medianos típicos de constelaciones LEO.

\subsection{6.2 Recomendaciones para Implementación}

Uno de los aspectos más prácticos que emerge de esta investigación radica en las condiciones necesarias para trasladar estos resultados desde el laboratorio a la línea de producción. La orientación de las piezas durante la impresión, el control exhaustivo de parámetros de proceso y una estrategia inteligente de post-procesado (especialmente en superficies RF) resultan decisivos para reproducir consistentemente los beneficios observados.

Se recomienda comenzar implementaciones reales con componentes secundarios o brackets no primarios, donde el riesgo es menor, para acumular experiencia de calificación antes de abordar estructuras principales. Asimismo, la integración temprana de equipos de aseguramiento de calidad en la fase de DfAM evita iteraciones costosas downstream. Aunque la tecnología ya muestra madurez suficiente, ignorar la curva de aprendizaje en metrología y control de distorsión térmica puede diluir gran parte de las ventajas prometidas.

\subsection{6.3 Trabajos Futuros}

Lejos de considerarse un capítulo cerrado, este campo presenta numerosas avenidas prometedoras. Dalmarco et al. \cite{Dalmarco2026} anticipan un crecimiento sostenido en la adopción de AM en el sector espacial, condicionado principalmente a avances en calificación y estandarización. En esta línea, resulta prioritario desarrollar bases de datos certificadas de propiedades de materiales impresos bajo condiciones espaciales reales, incluyendo degradación por radiación y ciclos térmicos extremos.

Entre las direcciones más interesantes destacan la exploración de manufactura híbrida (combinando aditivo con procesos convencionales selectivos), el desarrollo de materiales multifuncionales con propiedades RF y estructurales optimizadas, y la escalabilidad hacia impresión en órbita. Esta última, aunque todavía incipiente, podría revolucionar completamente los paradigmas de mantenimiento y extensión de vida útil de constelaciones.

Otros frentes relevantes incluyen la optimización multiobjetivo bajo incertidumbre (considerando variabilidad de proceso) y el análisis de ciclo de vida completo desde una perspectiva de sostenibilidad. Avanzar en estas direcciones no solo fortalecerá la competitividad técnica, sino que contribuirá a hacer más accesible y sostenible el despliegue de infraestructuras de telecomunicaciones en el espacio.

En definitiva, la fabricación aditiva ha dejado de ser una tecnología emergente para convertirse en una herramienta estratégica. Su implementación inteligente, acompañada de esfuerzos concertados en calificación y maduración, promete transformar la próxima generación de sistemas satelitales de telecomunicaciones.